\subsection{Performance Comparison}

\begin{table*}[t]
\centering
\renewcommand{\arraystretch}{1.2} % 调整行间距
\setlength{\tabcolsep}{9pt} % 调整列间距为12pt（默认是6pt）
\caption{Performance comparison over all methods. The bold and underlined values indicate the best and runner-up results, respectively.}
% \large
% \resizebox{\textwidth}{!}{
\begin{tabular}{lcccccccc}
\hline
          & \multicolumn{2}{c}{\textbf{Tweet}} & \multicolumn{2}{c}{\textbf{News-T}} & \multicolumn{2}{c}{\textbf{News-S}} & \multicolumn{2}{c}{\textbf{News-TS}} \\ \cline{2-9} 
                         & \textbf{ACC} & \textbf{NMI}       & \textbf{ACC} & \textbf{NMI}       & \textbf{ACC} & \textbf{NMI}       & \textbf{ACC} & \textbf{NMI}        \\ \hline
K-means(TF-IDF)          & 0.563        & 0.795              & 0.562        & 0.777              & 0.621        & 0.831              & 0.684        & 0.879               \\
K-means(Embedding)       & 0.581        & 0.839              & 0.589        & 0.806              & 0.613        & 0.831              & 0.665        & 0.876               \\
HieClu(TF-IDF)          & 0.589        & 0.784              & 0.593        & 0.773               & 0.702        & 0.845              & 0.788        & 0.901               \\
K-means(Embedding)       & 0.628        & 0.857              & 0.641        & 0.823               & 0.681        & 0.845              & 0.756        & 0.895               \\
LDA                      & 0.609        & 0.805              & 0.680        & 0.824              & 0.769        & 0.871              & 0.801        & 0.908               \\
BTM                      & 0.575        & 0.807              & 0.715        & 0.875              & 0.738        & 0.891              & 0.762        & 0.920               \\
WNTM                     & 0.678        & 0.851              & 0.781        & 0.878              & 0.743        & 0.863              & 0.814        & 0.918               \\
PTM                      & 0.685        & 0.840              & 0.738        & 0.866              & 0.785        & 0.878              & 0.814        & 0.911               \\
% GPU-DMM                  & 0.532        & 0.809              & 0.654        & 0.849              & 0.742        & 0.874              & 0.774        & 0.905               \\
% GPU-PDMM                 & 0.610        & 0.827              & 0.700        & 0.870              & 0.786        & 0.886              & 0.840        & 0.926               \\
% LFDMM                    & 0.607        & 0.845              & 0.736        & 0.874              & 0.766        & 0.887              & 0.800        & 0.920               \\
LCTM                     & 0.614        & 0.797              & 0.636        & 0.784              & 0.726        & 0.863              & 0.791        & 0.894               \\
SCCL                     & 0.782        & 0.892              & 0.758        & 0.883              & 0.831        & 0.904              & \textbf{0.898}        & \underline{0.949}               \\ 
DACL                     & \underline{0.820}         & 0.906              &0.805        & 0.898               & 0.840        & \underline{0.917}              & \underline{0.896}         & 0.944               \\ 
RSTC                     & 0.752        & 0.873              & 0.722        & 0.873               & 0.793        & 0.894              & 0.832        & 0.931               \\ 
MVC                      &0.807         & \underline{0.910}              & \textbf{0.847}        & \textbf{0.918}     
        & \underline{0.845}        & 0.916              & 0.878        & 0.946               \\ \hline
GSDMM                    & 0.778        & 0.895             & 0.778        & 0.892              & 0.797       & 0.901              & 0.829        & 0.934               \\
\sysname{}   & \textbf{0.870} & \textbf{0.914}     & \underline{0.830} & \underline{0.906}    & \textbf{0.855} & \textbf{0.929}     & 0.889 & \textbf{0.957}      \\ \hline
\end{tabular}
% }
\label{tab:performance_comparison}
\end{table*}

Table~\ref{tab:performance_comparison} compares the proposed GSDMM and \sysname{} against a range of baseline methods across four datasets: Tweet, News-T, News-S, and News-TS. Accuracy (ACC) and Normalized Mutual Information (NMI) are the evaluation metrics employed, which collectively assess the clustering performance from both classification and information-theoretic perspectives.

From the table, we can infer the following observations: 
1) GSDMM outperforms most methods while maintaining simplicity, elegance, and high efficiency. With enhancements, \sysname{} achieves significant performance improvements across all aspects.
2) \sysname{} consistently outperforms all baseline methods regarding ACC and NMI on the Tweet and News-S datasets. Notably, \sysname{} achieves the best clustering performance without leveraging text embeddings, demonstrating its superiority. In particular, we improve accuracy by 5\% on the Tweet dataset.
3) Compared to MVC, which performs text clustering using both the bag-of-words and text embedding views, our methods relies solely on the bag-of-words representation. This design ensures extremely high efficiency while still achieving competitive performance. For the News-T dataset, \sysname{} achieves the second-best performance, with an ACC of 0.83 and an NMI of 0.906, while maintaining high efficiency. 
4) \sysname{} demonstrates superior performance compared to SCCL, which uses the sentence-bert model and fine-tunes it with contrastive loss and clustering loss. The method exploits word co-occurrences in texts and is highly competitive compared to embedding-based models, thereby providing a significant advancement over existing methods.
